\bta{2012}


\begin{enumerate}
	%\renewcommand{\labelenumi}{\arabic{enumi}.}
	% A(\Alph) a(\alph) I(\Roman) i(\roman) 1(\arabic)
	%设定全局标号series=example	%引用全局变量resume=example
	%[topsep=-0.3em,parsep=-0.3em,itemsep=-0.3em,partopsep=-0.3em]
	%可使用leftmargin调整列表环境左边的空白长度 [leftmargin=0em]
	\item
	选择题 (本题满分 50 分，每小题 5 分)
	\begin{enumerate}
		%\renewcommand{\labelenumi}{\arabic{enumi}.}
		% A(\Alph) a(\alph) I(\Roman) i(\roman) 1(\arabic)
		%设定全局标号series=example	%引用全局变量resume=example
		%[topsep=-0.3em,parsep=-0.3em,itemsep=-0.3em,partopsep=-0.3em]
		%可使用leftmargin调整列表环境左边的空白长度 [leftmargin=0em]
		\item
		函数 $f(x)=x \cos x^{2}$, 正确结论是 \xzanswer{C} 
		
	\fourchoices
	{在 $(-\infty, 0)$ 内有界}
	{当 $x \rightarrow \infty$ 时, $f(x)$ 为无穷大}
	{在 $(-\infty, 0)$ 内无界}
	{当 $x \rightarrow \infty$ 时, $f(x)$ 极限存在}
	
\item 
函数 $f(x)$ 在 $(-\infty,+\infty)$ 上是连续函数, 且 $0<m<f(x)<M<\infty$, 则 $\frac{1}{m} \int_{-m}^{m}(f(t)-M) d t$ 的最大取值区间是 \xzanswer{B}


\fourchoices
{$(-M-m, m-M)$}
{$(2 m-2 M, 0)$}
{$(m-M, 0)$}
{$(0, M+m)$}


\item 
微分方程 $y y^{\prime \prime}-\left(y^{\prime}\right)^{2}=0$ 的一个特解是 \xzanswer{D} 


\fourchoices
{$y=x e^{x}$}
{$y=x \ln x$}
{$y=\ln x$}
{$y=e^{x}$}

\item 
已知 $n, m$ 是正整数, 且 $n<m$, 如果 $A=\int_{0}^{1} x^{m}(1-x)^{n} d x, B=\int_{0}^{1} x^{n}(1-x)^{m+1} d x$, 则下面结论正确的一个是 \xzanswer{A} 


\fourchoices
{$A>B$}
{$A=B$}
{$A<B$}
{$A, B$ 的大小关系不确定}

\item 
函数 $f(x)=e^{x}-x^{2}-4 x-3$ 在其定义域内零点的个数是 \xzanswer{C} 



\fourchoices
{$ 1 $}
{$ 2 $}
{$ 3 $}
{多于$  3 $}

\item 
若函数 $f(x)= \left\{\begin{aligned} &e^{x}(\sin x+\cos x) &   x \geqslant 0 \\ &a b x^{2}+a x+2 a+b &   x<0\end{aligned}\right.$ 的导函数在 $(-\infty,+\infty)$ 上连续, 则 \xzanswer{B} 


\fourchoices
{$a=2, b=-1$}
{$a=2, b=-3$}
{$a=1, b=-3$}
{$a=1, b=-1$}

\item 
若幂级数 $\sum\limits_{n=1}^{\infty} a_{n}(x-1)^{n}$ 在 $x=4$ 处条件收敛, 则级数 $\sum\limits_{n=1}^{\infty}(-1)^{n}\left(1+2^{n}\right) a_{n}$ \xzanswer{C} 


\fourchoices
{条件收敛}
{发散}
{绝对收敛}
{不能确定}

\item 
设 $S$ 为螺旋面 $\left\{\begin{array}{l}x=u \cos v \\ y=u \sin v   \\ z=v\end{array} \right.$的一部份，$ 0\leq u \leq \sqrt{15} $，$ 0 \leq v \leq \pi $ 则 $\iint_{S} \sqrt{x^{2}+y^{2}} d S$的值为 \xzanswer{C} 


\fourchoices
{$17 \pi$}
{$19 \pi$}
{$21 \pi$}
{$23 \pi$}

\item 
$\lim\limits _{x \rightarrow 0}\left(\frac{x}{\sin x}\right)^{\frac{1}{1-\cos x}}$ 的值为 \xzanswer{A} 


\fourchoices
{$e^{\frac{1}{3}}$}
{$e^{\frac{-1}{3}}$}
{$e^{\frac{1}{2}}$}
{$e^{-\frac{1}{2}}$}


\item 
一平面过点 $M(1,1,-1)$ 且与直线 $L: \frac{x}{2}=\frac{y+1}{1}=\frac{z-3}{-1}$ 垂直, 则该平面与平 面 $x-2 y-z+1=0$ 的交线的方向数是 \xzanswer{D} 


\fourchoices
{$(-5,1,3)$}
{$(1,-3,5)$}
{$(1,-5,3)$}
{$(3,-1,5)$}




	\end{enumerate}


\banswer{
\begin{enumerate}
	%\renewcommand{\labelenumi}{\arabic{enumi}.}
	\item
	\textbf{C}。取 $x_k=-\sqrt{2k\pi}\to-\infty$，则 $\cos x_k^2=1$，
	$f\left(x_k\right)=-\sqrt{2k\pi}\to-\infty$，故 $f$ 在 $(-\infty,0)$ 内无界，C 正确、A 错。
	又取 $y_k=-\sqrt{\left(2k+\frac12\right)\pi}$ 得 $f\left(y_k\right)=0$，
	故 $x\to\infty$ 时 $f$ 既不是无穷大也无极限，B、D 皆错。

	\item
	\textbf{B}。由 $0<m<f(t)<M$ 得 $m-M<f(t)-M<0$，在长为 $2m$ 的区间上积分：
	\[
	2m(m-M)<\int_{-m}^{m}\left(f(t)-M\right)dt<0 .
	\]
	除以 $m>0$ 即得取值范围 $(2m-2M,\ 0)$，且这两个界可被任意逼近。

	\item
	\textbf{D}。逐一代入。$y=e^{x}$：$yy''-\left(y'\right)^2=e^{2x}-e^{2x}=0$，成立。
	其余三个：$y=xe^{x}$ 得 $-e^{2x}$；$y=x\ln x$ 得 $\ln x-(\ln x+1)^2$；
	$y=\ln x$ 得 $-\frac{\ln x+1}{x^2}$，均不恒为零。
	（事实上方程即 $\left(\frac{y'}{y}\right)'=0$，通解为 $y=C_1e^{C_2x}$。）

	\item
	\textbf{A}。用 Beta 函数 $\int_0^1x^{p}(1-x)^{q}dx=\frac{p!\,q!}{(p+q+1)!}$：
	\[
	A=\frac{m!\,n!}{(m+n+1)!},\qquad B=\frac{n!\,(m+1)!}{(m+n+2)!} .
	\]
	于是
	\[
	\frac{A}{B}=\frac{m!}{(m+1)!}\cdot\frac{(m+n+2)!}{(m+n+1)!}=\frac{m+n+2}{m+1}>1
	\]
	（因 $n\geqslant1$），故 $A>B$。（此结论与 $n<m$ 无关。）

	\item
	\textbf{C}。$f'(x)=e^{x}-2x-4$，$f''(x)=e^{x}-2$，故 $f'$ 在 $x=\ln2$ 处取最小值
	$f'(\ln2)=-2-2\ln2<0$；而 $x\to\pm\infty$ 时 $f'\to+\infty$，
	故 $f'$ 恰有两个零点，$f$ 先增、后减、再增，至多三个零点。
	又
	\[
	f(-4)=e^{-4}-3<0,\quad f(-3)=e^{-3}>0,\quad f(0)=-2<0,\quad f(4)>0,
	\]
	由介值定理在 $(-4,-3)$、$(-3,0)$、$(0,4)$ 内各有一个零点，故恰为 $3$ 个。
	（数值解约为 $-3.024$、$-0.795$、$3.299$。）

	\item
	\textbf{B}。连续性：$f\left(0^{+}\right)=1$，$f\left(0^{-}\right)=2a+b$，故 $2a+b=1$。
	导数：$x>0$ 时 $f'=2e^{x}\cos x\to2$；$x<0$ 时 $f'=2abx+a\to a$。
	$f'$ 连续要求 $a=2$，再由 $2a+b=1$ 得 $b=-3$。

	\item
	\textbf{C}。在 $x=4$ 条件收敛说明 $|x-1|=3$ 恰为收敛半径，即 $R=3$。
	所求级数
	\[
	\sum(-1)^n\left(1+2^n\right)a_n=\sum a_n(-1)^n+\sum a_n(-2)^n ,
	\]
	两个级数分别对应 $x-1=-1$ 与 $x-1=-2$，
	而 $|-1|<3$、$|-2|<3$ 都在收敛区间内部，故都绝对收敛，其和亦绝对收敛。

	\item
	\textbf{C}。$\bm{r}_u=(\cos v,\sin v,0)$，$\bm{r}_v=(-u\sin v,u\cos v,1)$，
	\[
	\bm{r}_u\times\bm{r}_v=(\sin v,-\cos v,u),\quad
	\left|\bm{r}_u\times\bm{r}_v\right|=\sqrt{1+u^2}.
	\]
	又 $\sqrt{x^2+y^2}=u$，故
	\[
	\iint\limits_{S}\sqrt{x^2+y^2}\,dS=\int_0^{\pi}\!\!dv\int_0^{\sqrt{15}}u\sqrt{1+u^2}\,du
	=\pi\cdot\frac{16^{3/2}-1}{3}=21\pi .
	\]

	\item
	\textbf{A}。由 $\frac{x}{\sin x}=\left(1-\frac{x^2}{6}+o\left(x^2\right)\right)^{-1}
	=1+\frac{x^2}{6}+o\left(x^2\right)$ 得 $\ln\frac{x}{\sin x}\sim\frac{x^2}{6}$，
	而 $1-\cos x\sim\frac{x^2}{2}$，故指数
	$\frac{\ln\frac{x}{\sin x}}{1-\cos x}\to\frac13$，原式 $\to e^{1/3}$。

	\item
	\textbf{D}。所求平面以 $L$ 的方向 $(2,1,-1)$ 为法向量，过 $M(1,1,-1)$，
	即 $2x+y-z-4=0$。它与 $x-2y-z+1=0$ 的交线方向为
	\[
	(2,1,-1)\times(1,-2,-1)=(-3,1,-5)\parallel(3,-1,5) .
	\]
\end{enumerate}
}


\item
(本题满分 10 分)
证明极限 $\lim\limits _{n \rightarrow \infty}\left(\frac{1}{n}+\frac{1}{n+1}+\cdots+\frac{1}{3 n}\right)$ 存在, 并求出极限值。

\banswer{
	记 $S_n=\sum\limits_{k=n}^{3n}\frac1k$。

	\textbf{极限存在.}
	一方面 $S_n>0$ 有下界；另一方面
	\begin{align*}
	S_{n+1}-S_n&=\frac{1}{3n+1}+\frac{1}{3n+2}+\frac{1}{3n+3}-\frac1n\\
	&<\frac{3}{3n+1}-\frac1n=\frac{-1}{n(3n+1)}<0 ,
	\end{align*}
	故 $\left\{S_n\right\}$ 单调递减有下界，由单调有界定理极限存在。

	\textbf{求极限（定积分定义）.}
	把 $S_n$ 写成 Riemann 和：令 $k=n+j\ (j=0,1,\dots,2n)$，
	\[
	S_n=\sum_{j=0}^{2n}\frac{1}{n+j}=\frac1n\sum_{j=0}^{2n}\frac{1}{1+\frac{j}{n}} .
	\]
	右端是函数 $\frac{1}{1+x}$ 在 $[0,2]$ 上以 $\frac1n$ 为步长的 Riemann 和
	（多出的一个端点项 $\frac1n\cdot\frac13\to0$ 不影响极限），故
	\[
	\lim_{n\to\infty}S_n=\int_0^2\frac{dx}{1+x}=\ln3 .
	\]
	（数值验证：$n=10^6$ 时 $S_n\approx1.0986130$，而 $\ln3\approx1.0986123$。）
}


\item 
(本题满分 10 分)	
求微分方程 $y^{\prime \prime}-3 y^{\prime}+2 y=e^{x}(2 x+1)$ 的通解。

\banswer{
	\textbf{齐次通解.}
	特征方程 $r^2-3r+2=0$ 的根为 $r=1,2$，故
	$Y=C_1e^{x}+C_2e^{2x}$。

	\textbf{特解.}
	右端为 $e^{\lambda x}P_1(x)$，$\lambda=1$ 恰是单重特征根，故设
	\[
	y^{*}=x\left(Ax+B\right)e^{x}=\left(Ax^2+Bx\right)e^{x}.
	\]
	求导：
	\begin{align*}
	y^{*\prime}&=\left(Ax^2+(2A+B)x+B\right)e^{x},\\
	y^{*\prime\prime}&=\left(Ax^2+(4A+B)x+2A+2B\right)e^{x}.
	\end{align*}
	代入左端，$x^2$ 项自动抵消，得
	\[
	y^{*\prime\prime}-3y^{*\prime}+2y^{*}=\left(-2Ax+2A-B\right)e^{x}.
	\]
	与 $e^{x}(2x+1)$ 比较：$-2A=2$、$2A-B=1$，即 $A=-1$、$B=-3$。

	\textbf{通解.}
	\[
	y=C_1e^{x}+C_2e^{2x}-\left(x^2+3x\right)e^{x}.
	\]
	（代回原方程符号验算无误。）
}



\item 
(本题满分 10 分)	
计算 $\iint_{D}(x|y|+x y) d x d y$, 其中 $D$ 是由抛物线 $5 y=x^{2}-6$ 和抛物线 $y^{2}=x$ 围成的 闭区域。


\banswer{
	\textbf{定区域.}
	由 $x=y^2$ 代入 $5y=x^2-6$ 得 $y^4-5y-6=0$，因式分解
	\[
	(y+1)(y-2)\left(y^2+y+3\right)=0 ,
	\]
	末项判别式 $1-12<0$ 无实根，故交点对应 $y=-1$ 与 $y=2$，
	即 $(1,-1)$ 与 $(4,2)$。对固定的 $y\in[-1,2]$，区域是
	\[
	D:\ -1\leqslant y\leqslant2,\quad y^2\leqslant x\leqslant\sqrt{5y+6}
	\]
	（左边界是抛物线 $x=y^2$，右边界是 $5y=x^2-6$ 的 $x>0$ 分支）。

	\textbf{利用被积函数的分段性.}
	$x|y|+xy=x\left(|y|+y\right)$，而 $|y|+y=2y\ (y>0)$、$=0\ (y\leqslant0)$。
	故 $y\leqslant0$ 的部分贡献为零，只需算 $0\leqslant y\leqslant2$：
	\begin{align*}
	I&=2\int_0^2y\,dy\int_{y^2}^{\sqrt{5y+6}}x\,dx
	=\int_0^2y\left[\left(5y+6\right)-y^4\right]dy\\
	&=\int_0^2\left(5y^2+6y-y^5\right)dy
	=\left[\frac{5y^3}{3}+3y^2-\frac{y^6}{6}\right]_0^2 .
	\end{align*}
	代入得
	\[
	I=\frac{40}{3}+12-\frac{32}{3}=\frac{8}{3}+12=\frac{44}{3}.
	\]
	（数值积分同样给出 $14.6667$。）
}


\item 
(本题满分 10 分)	
将函数 $f(x)=|x-1|(0 \leqslant x \leqslant \pi)$ 展开成正弦级数。


\banswer{
	\textbf{作奇延拓.}
	把 $f$ 奇延拓到 $[-\pi,\pi]$ 再作 $2\pi$ 周期延拓，则 $a_n=0$，
	\[
	b_n=\frac{2}{\pi}\int_0^{\pi}\left|x-1\right|\sin nx\,dx .
	\]

	\textbf{分段积分.}
	以 $x=1$ 为界拆开。由分部积分，
	\begin{align*}
	\int_0^{1}(1-x)\sin nx\,dx&=\frac1n-\frac{\sin n}{n^2},\\
	\int_1^{\pi}(x-1)\sin nx\,dx&=\frac{(1-\pi)(-1)^{n}}{n}-\frac{\sin n}{n^2}.
	\end{align*}
	两式相加再乘 $\frac{2}{\pi}$：
	\[
	b_n=\frac{2}{\pi}\left[\frac1n+\frac{(-1)^{n}(1-\pi)}{n}-\frac{2\sin n}{n^2}\right].
	\]

	\textbf{结论.}
	\[
	\left|x-1\right|=\sum_{n=1}^{\infty}
	\frac{2}{\pi}\left[\frac{1+(-1)^{n}(1-\pi)}{n}-\frac{2\sin n}{n^2}\right]\sin nx
	\]
	对 $0<x<\pi$ 成立（延拓后的函数在 $(0,\pi)$ 内连续且分段光滑）。
	在 $x=0$ 与 $x=\pi$ 处，正弦级数的每一项都为 $0$，故都收敛于 $0$；
	这与 Dirichlet 定理一致：奇延拓在 $x=0$ 处左右极限为 $\mp1$、
	在 $x=\pi$ 处为 $\pm(\pi-1)$，平均值均为 $0$。

	\par\vspace{1mm}\noindent
	\textbf{注：}{\kaishu
	原答案把系数中的 $(-1)^{n}$ 误排成 $(-1)^{2}$（即恒为 $1$），应予改正。
	按改正后的系数求部分和，在 $x=0.3,1,2,3$ 处分别得
	$0.6999,\ 0.0000,\ 1.0000,\ 2.0001$，与 $|x-1|$ 吻合。
	}\par
}


\item 
(本题满分 10 分)	
设函数 $f(x)=\int_{x}^{1} e^{-t^{2}} d t$, 求 $\int_{0}^{1} x^{2} f(x) d x$ 的值。

\banswer{
	\textbf{解法一（分部积分）.}
	注意 $f(1)=0$，$f'(x)=-e^{-x^2}$。取 $u=f(x)$、$dv=x^2dx$：
	\begin{align*}
	\int_0^1x^2f(x)\,dx&=\left[\frac{x^3}{3}f(x)\right]_0^1-\frac13\int_0^1x^3f'(x)\,dx\\
	&=0+\frac13\int_0^1x^3e^{-x^2}dx .
	\end{align*}
	再令 $u=x^2$（$du=2x\,dx$）：
	\[
	\int_0^1x^3e^{-x^2}dx=\frac12\int_0^1ue^{-u}du
	=\frac12\left[-(u+1)e^{-u}\right]_0^1=\frac12\left(1-\frac2e\right).
	\]
	故原式 $=\frac16\left(1-\frac2e\right)$。

	\textbf{解法二（交换积分次序）.}
	\[
	\int_0^1x^2\!\left(\int_x^1e^{-t^2}dt\right)dx
	=\int_0^1e^{-t^2}\!\left(\int_0^{t}x^2dx\right)dt
	=\frac13\int_0^1t^3e^{-t^2}dt ,
	\]
	与解法一殊途同归。

	\[
	\int_0^1x^2f(x)\,dx=\frac{1}{6}\left(1-\frac{2}{e}\right)\approx0.04404 .
	\]
	（数值积分给出 $0.0440402$，一致。）
}


\newpage
\item 
(本题满分 10 分)	
计算曲线积分 $I=\oint_{L} \frac{x d y-y d x}{x^{2}+2y^{2}}$, 其中 $L$ 是由直线 $x+y=1, y=x-1$ 和半圆周 $x^{2}+y^{2}=1, x \leqslant 0$ 所围成的闭曲线, 方向为逆时针方向。


\banswer{
	记 $P=\frac{-y}{x^2+2y^2}$，$Q=\frac{x}{x^2+2y^2}$。

	\textbf{闭形式.}
	在 $(x,y)\neq(0,0)$ 处
	\[
	\frac{\partial Q}{\partial x}=\frac{2y^2-x^2}{\left(x^2+2y^2\right)^2}
	=\frac{\partial P}{\partial y},
	\]
	故在不含原点的单连通区域内积分与路径无关。但 $L$ 内部含奇点 $(0,0)$
	（$L$ 依次由 $x\leqslant0$ 的半圆 $(0,1)\to(-1,0)\to(0,-1)$、
	直线 $y=x-1$ 的 $(0,-1)\to(1,0)$、直线 $x+y=1$ 的 $(1,0)\to(0,1)$ 组成，
	正向且围住原点），不能直接用格林公式。

	\textbf{挖去奇点.}
	取充分小的 $\varepsilon>0$，作椭圆
	$\ell_\varepsilon:\ x^2+2y^2=\varepsilon^2$，取逆时针方向，
	使它完全落在 $L$ 内部。在 $L$ 与 $\ell_\varepsilon$ 所夹的复连通区域上用格林公式，
	被积式为零，故
	\[
	I=\oint_{L}=\oint_{\ell_\varepsilon}\frac{x\,dy-y\,dx}{x^2+2y^2}
	=\frac{1}{\varepsilon^2}\oint_{\ell_\varepsilon}\left(x\,dy-y\,dx\right),
	\]
	这里用了在 $\ell_\varepsilon$ 上分母恒为 $\varepsilon^2$。

	\textbf{算出.}
	由格林公式 $\oint(x\,dy-y\,dx)=2\times$ 所围面积。
	椭圆 $x^2+2y^2=\varepsilon^2$ 的半轴为 $\varepsilon$ 与 $\frac{\varepsilon}{\sqrt2}$，
	面积 $=\pi\cdot\varepsilon\cdot\frac{\varepsilon}{\sqrt2}=\frac{\pi\varepsilon^2}{\sqrt2}$。故
	\[
	I=\frac{1}{\varepsilon^2}\cdot\frac{2\pi\varepsilon^2}{\sqrt2}=\sqrt2\,\pi .
	\]
	（沿原闭曲线作数值线积分得 $4.4429$，而 $\sqrt2\pi\approx4.4429$，一致。）
}


\item 
(本题满分 10 分)	
设函数 $f(x)$ 连续, 且 $f^{2}(x) \leqslant|x|^{3}$, 记 $F(x)=\int_{0}^{1} f(x t) d t$, 求 $F^{\prime}(x)$, 并讨论 $F^{\prime}(x)$ 的连续性。

\banswer{
	\textbf{预备.}
	由 $f^2(x)\leqslant|x|^3$ 得 $|f(x)|\leqslant|x|^{3/2}$，特别地 $f(0)=0$，
	并且
	\[
	\left|\int_0^{x}f(u)\,du\right|\leqslant\int_0^{|x|}u^{3/2}du=\frac25|x|^{5/2}.
	\]

	\textbf{$x\neq0$ 时.}
	作代换 $u=xt$ 得 $F(x)=\frac1x\int_0^{x}f(u)\,du$，故
	\[
	F'(x)=-\frac{1}{x^2}\int_0^{x}f(u)\,du+\frac{f(x)}{x}.
	\]

	\textbf{$x=0$ 时.}
	$F(0)=\int_0^1f(0)\,dt=0$，用导数定义：
	\[
	\left|\frac{F(x)-F(0)}{x}\right|=\left|\frac{1}{x^2}\int_0^{x}f(u)du\right|
	\leqslant\frac25|x|^{1/2}\xrightarrow[x\to0]{}0 ,
	\]
	故 $F'(0)=0$。合起来
	\[
	F^{\prime}(x)=\left\{\begin{aligned}
	&-\frac{1}{x^{2}}\int_{0}^{x}f(u)\,du+\frac{f(x)}{x}, &&x\neq0,\\
	&0, &&x=0 .\end{aligned}\right.
	\]

	\textbf{连续性.}
	$x\neq0$ 处 $F'$ 由连续函数复合而成，连续。在 $x=0$ 处，
	\[
	\left|F'(x)\right|\leqslant\frac25|x|^{1/2}+|x|^{1/2}=\frac75|x|^{1/2}\to0=F'(0),
	\]
	故 $F'$ 在 $x=0$ 也连续。所以 $F'(x)$ 在整个实轴上连续。
}



\item 
(本题满分 10 分)	
函数 $f(x)$ 在 $[a, b]$ 上连续, 在 $(a, b)$ 内可导, 并且有 $0<a<b$ 。 证明:
存在 $\xi \in(a, b)$, 使得 $\frac{a+b}{2 \xi} f^{\prime}(\xi)=\frac{f(b)-f(a)}{b-a}$。


\banswer{
	取辅助函数 $g(x)=x^2$。由 $0<a<b$ 知 $g'(x)=2x\neq0$ 于 $(a,b)$，
	且 $f,g$ 在 $[a,b]$ 上连续、$(a,b)$ 内可导。由 Cauchy 中值定理，存在 $\xi\in(a,b)$ 使
	\[
	\frac{f(b)-f(a)}{b^2-a^2}=\frac{f'(\xi)}{2\xi}.
	\]
	而 $b^2-a^2=(b-a)(b+a)$，于是
	\[
	\frac{f(b)-f(a)}{b-a}=\frac{(a+b)f'(\xi)}{2\xi}=\frac{a+b}{2\xi}f'(\xi),
	\]
	即为所证。

	\par\vspace{1mm}\noindent
	\textbf{注：}{\kaishu
	条件 $0<a<b$ 是为了保证 $g'(x)=2x$ 在 $(a,b)$ 上不为零（否则 Cauchy 中值定理失效），
	同时保证结论中的 $\frac{a+b}{2\xi}$ 有意义。
	}\par
}


\item 
(本题满分 10 分)	
函数 $f(x)$ 在 $(-\infty,+\infty)$ 上满足 $f^{\prime \prime}(x)>0$ 。 证明:
\[
f\left(\frac{x_{1}+x_{2}+\cdots+x_{n}}{n}\right) \leq \frac{f\left(x_{1}\right)+f\left(x_{2}\right)+\cdots+f\left(x_{n}\right)}{n}
\]


\banswer{
	记 $\bar{x}=\frac{x_1+x_2+\cdots+x_n}{n}$。

	\textbf{在 $\bar{x}$ 处二阶展开.}
	对每个 $i$，由带 Lagrange 余项的 Taylor 公式，存在 $\xi_i$ 介于 $x_i$ 与 $\bar{x}$ 之间，
	\[
	f\left(x_i\right)=f(\bar{x})+f'(\bar{x})\left(x_i-\bar{x}\right)
	+\frac{f''\left(\xi_i\right)}{2}\left(x_i-\bar{x}\right)^2 .
	\]
	由 $f''>0$，末项非负，故
	\[
	f\left(x_i\right)\geqslant f(\bar{x})+f'(\bar{x})\left(x_i-\bar{x}\right).
	\]

	\textbf{求和.}
	对 $i=1,\dots,n$ 相加，注意 $\sum\limits_{i=1}^{n}\left(x_i-\bar{x}\right)=0$：
	\[
	\sum_{i=1}^{n}f\left(x_i\right)\geqslant nf(\bar{x})+f'(\bar{x})\cdot0=nf(\bar{x}).
	\]
	两端除以 $n$ 即得
	\[
	f\left(\frac{x_1+x_2+\cdots+x_n}{n}\right)\leqslant\frac{f\left(x_1\right)+\cdots+f\left(x_n\right)}{n}.
	\]
	等号成立当且仅当 $x_1=x_2=\cdots=x_n$。
}



\item 
(本题满分 10 分)	
设 $a<b$, 函数 $f(x)$ 在 $[a, b]$ 上连续, 且
\[
\int_{a}^{b} f(x) d x=\int_{a}^{b} x f(x) d x=\int_{a}^{b} x^{2} f(x) d x=0
\]
证明：在 $(a, b)$ 上至少存在三个不同点 $x_{1}, x_{2}, x_{3}$, 使得 $f\left(x_{1}\right)=f\left(x_{2}\right)=f\left(x_{3}\right)=0$。


\banswer{
	\textbf{解法一（三重原函数 + Rolle 定理）.}
	令
	\[
	F(x)=\int_a^{x}f(t)\,dt,\quad G(x)=\int_a^{x}F(t)\,dt,\quad H(x)=\int_a^{x}G(t)\,dt .
	\]
	显然 $F(a)=G(a)=H(a)=0$，且 $H'''=f$。

	先证 $F(b)=G(b)=H(b)=0$：
	\begin{itemize}
	\item $F(b)=\int_a^bf\,dx=0$（第一个条件）；
	\item 由分部积分
	$0=\int_a^bxf\,dx=\left[xF\right]_a^b-\int_a^bF\,dx=bF(b)-G(b)=-G(b)$，
	故 $G(b)=0$；
	\item 同理
	$0=\int_a^bx^2f\,dx=\left[x^2F\right]_a^b-2\int_a^bxF\,dx=-2\int_a^bxF\,dx$，
	而 $\int_a^bxF\,dx=\left[xG\right]_a^b-\int_a^bG\,dx=bG(b)-H(b)=-H(b)$，
	故 $H(b)=0$。
	\end{itemize}

	接着逐层用 Rolle 定理：
	\begin{itemize}
	\item $H(a)=H(b)=0\Rightarrow\exists\,c\in(a,b)$ 使 $H'(c)=G(c)=0$；
	\item $G$ 在 $a<c<b$ 三点取零 $\Rightarrow\exists\,d_1\in(a,c)$、$d_2\in(c,b)$ 使 $F=0$；
	\item $F$ 在 $a<d_1<d_2<b$ 四点取零 $\Rightarrow$ 在三个相邻区间内各得一点
	$x_1\in\left(a,d_1\right)$、$x_2\in\left(d_1,d_2\right)$、$x_3\in\left(d_2,b\right)$
	使 $F'=f=0$。
	\end{itemize}
	三点互异且都在 $(a,b)$ 内，证毕。

	\textbf{解法二（反证：变号次数）.}
	若 $f$ 在 $(a,b)$ 内变号的点不足三个，设为 $t_1<\cdots<t_k\ (k\leqslant2)$，
	取 $p(x)=\prod\limits_{i=1}^{k}\left(x-t_i\right)$，它是次数 $\leqslant2$ 的多项式，
	于是 $p(x)f(x)$ 在 $(a,b)$ 上不变号且不恒为零，故 $\int_a^bp(x)f(x)\,dx\neq0$；
	但由题设 $\int_a^bf=\int_a^bxf=\int_a^bx^2f=0$ 知该积分为 $0$，矛盾。
	故 $f$ 至少变号三次，从而至少有三个零点。
}




	
\end{enumerate}

